% Options for packages loaded elsewhere
\PassOptionsToPackage{unicode}{hyperref}
\PassOptionsToPackage{hyphens}{url}
%
\documentclass[
  11pt,
]{article}
\usepackage{amsmath,amssymb}
\usepackage{iftex}
\ifPDFTeX
  \usepackage[T1]{fontenc}
  \usepackage[utf8]{inputenc}
  \usepackage{textcomp} % provide euro and other symbols
\else % if luatex or xetex
  \usepackage{unicode-math} % this also loads fontspec
  \defaultfontfeatures{Scale=MatchLowercase}
  \defaultfontfeatures[\rmfamily]{Ligatures=TeX,Scale=1}
\fi
\usepackage{lmodern}
\ifPDFTeX\else
  % xetex/luatex font selection
\fi
% Use upquote if available, for straight quotes in verbatim environments
\IfFileExists{upquote.sty}{\usepackage{upquote}}{}
\IfFileExists{microtype.sty}{% use microtype if available
  \usepackage[]{microtype}
  \UseMicrotypeSet[protrusion]{basicmath} % disable protrusion for tt fonts
}{}
\makeatletter
\@ifundefined{KOMAClassName}{% if non-KOMA class
  \IfFileExists{parskip.sty}{%
    \usepackage{parskip}
  }{% else
    \setlength{\parindent}{0pt}
    \setlength{\parskip}{6pt plus 2pt minus 1pt}}
}{% if KOMA class
  \KOMAoptions{parskip=half}}
\makeatother
\usepackage{xcolor}
\usepackage[margin=23mm]{geometry}
\setlength{\emergencystretch}{3em} % prevent overfull lines
\providecommand{\tightlist}{%
  \setlength{\itemsep}{0pt}\setlength{\parskip}{0pt}}
\setcounter{secnumdepth}{5}
\allowdisplaybreaks
\setlength{\emergencystretch}{3em}
\ifLuaTeX
  \usepackage{selnolig}  % disable illegal ligatures
\fi
\usepackage{bookmark}
\IfFileExists{xurl.sty}{\usepackage{xurl}}{} % add URL line breaks if available
\urlstyle{same}
\hypersetup{
  pdftitle={Residue-driven constituent curvature and the original arithmetic Laplacian},
  pdfauthor={Research continuation},
  hidelinks,
  pdfcreator={LaTeX via pandoc}}

\title{Residue-driven constituent curvature and the original arithmetic
Laplacian}
\usepackage{etoolbox}
\makeatletter
\providecommand{\subtitle}[1]{% add subtitle to \maketitle
  \apptocmd{\@title}{\par {\large #1 \par}}{}{}
}
\makeatother
\subtitle{Exact rank-one transport, a source-sensitive curvature
coefficient, and the retained period correction}
\author{Research continuation}
\date{13 September 2026}

\begin{document}
\maketitle

\section{Scope and source
identifications}\label{scope-and-source-identifications}

This note continues the supplied \emph{Deligne Split Cohomology} sidebar
(DS1--69), \emph{Exponential Period Determinant} (XD1--26),
\emph{Exponential Translation Comparison} (DT1--13), and the pasted SGA
trace/period continuation. It uses their explicit coefficient and period
maps. The additional results proved here concern the rank-one
deformation on every nonzero proper invariant constituent, its exact
period curvature, the change of that coefficient between the original
source metrics, and the second-order realization of the arithmetic
Laplacian. These are written proofs, supplemented by the included exact
finite checker. They are not new Lean certificates or a proof of a
uniform arithmetic endpoint bound.

Keep the original nonempty cyclic object, including its full
multiplicities,

\[
E=\mathbb C[S]/(\chi),\qquad
\chi(S)=S^q+\sum_{a=0}^{q-1}c_aS^a,\qquad q\geq1,
\qquad A=M_S.
\tag{RC1}
\]

The coordinate basis is \(1,S,\ldots,S^{q-1}\). When \(E=E_{h,k}\) is
the arithmetic packet, the source retains \(g=2\xi\),
\(\mathcal MF_h=g/h\), the full arithmetic Taylor unit in \(\eta\), and
the original source/relation norm. In particular, retain the specified
maps on the admitted strong-Schwartz tensor domain \(\mathscr X\):

\[
r_N:E\longrightarrow\mathscr X,\quad j_E:\mathscr X\longrightarrow E,
\qquad j_Er_N=I_E,\quad j_ED^{(k)}=Aj_E,\quad j_Ed=0,
\qquad G_N=r_N^*r_N\succ0.
\tag{RC2}
\]

Here \(d\) is the original cochain differential and \(D^{(k)}\) the
original sum of scaling generators. Equation (RC2) is the supplied
source comparison, not an assertion that the full-jet observation
extends boundedly to the unweighted Hilbert completion (see
\emph{Exponential Intake}, §§4--5).

There are three distinct parameters. \(S\) is the arithmetic spectral
coordinate; \(z\) is the earlier representative-curve parameter; \(t\)
below is the exponential deformation parameter. The metric curve of
DS13--14 is \(G(z)=G_j+|z|^2(G_i-G_j)\). None of its radial derivatives
is substituted for an exponential \(t\)-derivative.

For fixed \(u\in\mathbb C^\times\), the supplied polynomial complex is

\[
\mathcal C_{u,t}=
[\mathbb C[S]\xrightarrow{\mathcal D_{u,t}}\mathbb C[S]dS],
\qquad \mathcal D_{u,t}P=(uP'+(\chi-t)P)dS.
\tag{RC3}
\]

Its degree-one cohomology is a coefficient-vector-space quotient with
the specified degree-below-\(q\) frame. Multiplication by \(S\) does not
descend as an algebra operation on this general fibre: the supplied
identities \([\chi-t]=0\) and \([S(\chi-t)]=-u[1]\) rule out that
interpretation. The parameter connection does descend, and in the
retained frame it is

\[
\nabla_t=\partial_t-\frac{A(t)}u,\qquad
A(t)=A+t\mathcal R,\qquad
\mathcal R=e_0e_{q-1}^{T}.
\tag{RC4}
\]

At \(u=t=0\) the complex specializes to the original quotient by
\(\chi\). The regular operator \(-u\nabla_t\) induces \(A\) on that
special fibre; this is the type-correct specialization in DS68.

The fixed oriented contours of XD2 give a period matrix
\(\Pi(t)=\Pi(u,\mathbf c,t)\) satisfying

\[
\Pi'(t)=-\frac1u\Pi(t)A(t),\qquad
\det\Pi(t)=\det\Pi_{\mathrm{mon}}(u)
\exp\left(\frac{\mathcal F(\mathbf c,t)}u\right),
\quad \mathcal F=\operatorname{Tr}\Phi_t(A(t)).
\tag{RC5}
\]

XD4--7 supply the uniform contour estimates that justify holomorphy and
differentiation. XD16--20 supply the nonzero monomial initial
determinant and an adjugate proof valid even at a hypothetical singular
matrix. Thus \(\Pi(t)\) is invertible for every complex \(t\), without
deleting repeated-root fibres. The gamma multiplication formula
simplifies XD18 to

\[
\det(\Pi(t)^*\Pi(t))=(2\pi|u|)^q
\exp\left(2\operatorname{Re}\frac{\mathcal F(\mathbf c,t)}u\right).
\tag{RC6}
\]

This last simplification is already in the supplied SGA continuation. It
is not a new result of the present note. Bloch and Esnault (1999, §5)
study period determinants of this polynomial-exponential type; the
literal normalizations used here are the ones proved in XD1--26.

\section{The deformation is the residue rank-one
map}\label{the-deformation-is-the-residue-rank-one-map}

Define the original residue functional and its bilinear pairing by

\[
\ell([P])=[S^{q-1}]\operatorname{rem}_\chi P,
\qquad \beta(x,y)=\ell(xy).
\tag{RC7}
\]

In the power basis, \(\beta(S^i,S^j)\) is zero for \(i+j<q-1\) and one
for \(i+j=q-1\). Reversing its columns therefore produces a triangular
matrix with diagonal one. Consequently

\[
\det\beta=(-1)^{q(q-1)/2},\qquad
\beta^\flat:E\xrightarrow{\sim}E^\vee.
\tag{RC8}
\]

This proof inverts no discriminant. It is the perfect residue pairing in
the supplied finite-flat SGA calculation. Its contraction with
multiplication by \(\chi'\) gives the separate, potentially degenerate
trace pairing:

\[
\operatorname{Tr}_E M_P=\ell(\chi'P).
\tag{RC9}
\]

For completeness, the divided difference
\(\mathcal C_\chi(X,Y)=(\chi(X)-\chi(Y))/(X-Y)\) is the coevaluation
tensor for (RC8): \((\ell\otimes1)((P\otimes1)\mathcal C_\chi)=P\) and
its image under multiplication is \(\chi'\). The finite-dimensional
evaluation--coevaluation trace then proves (RC9). Neither (RC8) nor
(RC9) identifies the bilinear residue form with the positive Hermitian
source form \(G_N\).

In the fixed remainder frame the perturbation in (RC4) is exactly

\[
\boxed{\mathcal R(v)=\ell(v)\,1_E,
\qquad \mathcal R=1_E\otimes\ell.}
\tag{RC10}
\]

It follows that the pair \((A,\mathcal R)\) generates every coefficient
endomorphism:

\[
\boxed{\mathbb C\langle A,\mathcal R\rangle
=\operatorname{End}_{\mathbb C}(E).}
\tag{RC11}
\]

\textbf{Proof.} The vectors \(A^i1_E\), \(0\leq i<q\), are the power
basis. By (RC8), the functionals \(\ell\circ A^j\), \(0\leq j<q\), are a
basis of \(E^\vee\). Therefore the \(q^2\) operators

\[
A^i\mathcal R A^j=(A^i1_E)\otimes(\ell\circ A^j)
\tag{RC12}
\]

are a vector-space basis of \(\operatorname{End}_{\mathbb C}(E)\). Every
one belongs to the algebra generated by \(A\) and \(\mathcal R\).

Now fix an actual nonzero proper invariant constituent, with its full
maps,

\[
0\longrightarrow F\xrightarrow{\iota}E
\xrightarrow{\pi}E/F\longrightarrow0,
\qquad A\iota=\iota A_F,
\qquad 0<p:=\dim F<q.
\tag{RC13}
\]

The deformation has a specified transverse map

\[
\boxed{\pi\mathcal R\iota
=(\pi1_E)\otimes\ell_F:F\longrightarrow E/F,
\qquad \ell_F=\ell\circ\iota.}
\tag{RC14}
\]

Both tensor factors are nonzero. If \(1_E\) belonged to \(F\), its
entire \(A\)-orbit would belong to \(F\), contradicting \(F\ne E\). If
\(\ell_F=0\), then for \(v\in F\), invariance would give
\(\ell(A^jv)=0\) for every \(j\). Nondegeneracy of (RC8) would force
every \(v\in F\) to be zero. Thus

\[
\boxed{\operatorname{rank}(\pi\mathcal R\iota)=1,
\quad \ker(\pi\mathcal R\iota)=\ker\ell_F,
\quad \operatorname{im}(\pi\mathcal R\iota)=\mathbb C\,\pi1_E.}
\tag{RC15}
\]

In particular the kernel has dimension \(p-1\). This is the
strengthening of the supplied rank-at-most-one statement. It holds with
all nilpotent jets retained. Equivalently, an invariant subspace is an
ideal \((a)/(\chi)\) for a monic divisor \(a\mid\chi\); the polynomial
\(S^{q-1-\deg a}a\) witnesses that \(\ell_F\) is nonzero whenever
\(0<\deg a<q\).

No nonzero proper \emph{constant remainder-coordinate} subspace is
invariant under \(A(t)\) for every \(t\). This assertion must not be
called differential irreducibility: moving horizontal subbundles do
exist. For example, \(\Pi(t)^{-1}\Pi(0)F\) is horizontal. It is a
different marked family from the fixed original constituent in (RC13).

\section{Exact constituent-period
curvature}\label{exact-constituent-period-curvature}

Keep \(F\) and its inclusion fixed as in (RC13), and put

\[
P_F(t)=\Pi(t)\iota,\qquad
H(t)=\Pi(t)^*\Pi(t),\qquad
H_F(t)=P_F(t)^*P_F(t)=\iota^*H(t)\iota.
\tag{RC16}
\]

The target metric in this section is exactly the fixed coordinate metric
\(I_q\) of the period construction. Since \(\Pi(t)\) is invertible,
\(P_F(t)\) has rank \(p\) everywhere. Its orthogonal projector is

\[
\mathsf P_F(t)=P_F(t)H_F(t)^{-1}P_F(t)^*.
\tag{RC17}
\]

Use (RC4) and \(A\iota=\iota A_F\) before projecting:

\[
P_F'(t)=-\frac1uP_F(t)A_F
-\frac t u\Pi(t)\mathcal R\iota.
\tag{RC18}
\]

The first term is tangent to the same period subspace. Thus its complete
normal derivative is

\[
\boxed{\mathcal N_F(t):=(I-\mathsf P_F(t))P_F'(t)
=-\frac t u (I-\mathsf P_F(t))\Pi(t)1_E\,\ell_F.}
\tag{RC19}
\]

The column \((I-\mathsf P_F(t))\Pi(t)1_E\) never vanishes: vanishing
would imply \(1_E\in F\). Equations (RC15) and (RC19) prove that this
normal map is zero at \(t=0\) and has rank exactly one at every
\(t\ne0\). Its kernel for \(t\ne0\) is \(\ker\ell_F\). At the marked
fibre,

\[
\boxed{(I-\mathsf P_F(0))P_F'(0)=0,
\qquad (I-\mathsf P_F(0))P_F''(0)
=-\frac1u(I-\mathsf P_F(0))\Pi(0)1_E\,\ell_F\ne0.}
\tag{RC20}
\]

The second identity follows by differentiating (RC5):
\(\Pi''=\Pi A(t)^2/u^2-\Pi\mathcal R/u\). At zero,
\(A^2\iota=\iota A_F^2\), so the \(A^2\) term vanishes under the stated
normal projection. The surviving map has rank one.

\subsection{The exact coefficient of
curvature}\label{the-exact-coefficient-of-curvature}

For any positive Hermitian matrix \(M\) on \(E\), define

\[
q_{M,F}([1_E])=\min_{f\in F}\|1_E-\iota f\|_M^2,
\qquad M_F=\iota^*M\iota,
\tag{RC21}
\]

and

\[
\boxed{\mathfrak c_F(M)
=q_{M,F}([1_E])\,\ell_F M_F^{-1}\ell_F^*.}
\tag{RC22}
\]

This scalar is strictly positive for (RC13). It is independent of the
basis chosen in \(F\), but it retains the marked quotient vector
\([1_E]\) and the marked residue functional \(\ell_F\). The two factors
are the actual quotient norm and the actual dual norm, not arbitrary
positive constants.

Let \(\psi_F(t)=\log\det H_F(t)\), where the determinant is of the
\(p\times p\) positive Gram. No rectangular determinant is used.
Ordinary matrix differentiation, using holomorphy of \(P_F\), gives

\[
\begin{split}
\partial_t\partial_{\bar t}\psi_F
&=\operatorname{Tr}\left[
H_F^{-1}P_F'^*P_F'
-H_F^{-1}P_F'^*P_FH_F^{-1}P_F^*P_F'
\right]\\
&=\operatorname{Tr}\left[H_F^{-1}\mathcal N_F^*\mathcal N_F\right].
\end{split}
\tag{RC23}
\]

The projector in (RC17) is self-adjoint and idempotent, which justifies
the last equality. Further,

\[
q_{H(t),F}([1_E])
=\|(I-\mathsf P_F(t))\Pi(t)1_E\|^2.
\tag{RC24}
\]

Substitute (RC19) and take the trace of its rank-one product. This
yields an exact identity for all complex \(t\):

\[
\boxed{\partial_t\partial_{\bar t}\log\det H_F(t)
=\frac{|t|^2}{|u|^2}\,\mathfrak c_F(H(t)).}
\tag{RC25}
\]

Thus the dual determinant line of the period-image subbundle, with its
specified induced metric, has Chern form

\[
c_1((\det\mathscr F_{\rm per})^\vee)
=\frac{i}{2\pi}\frac{|t|^2}{|u|^2}
\mathfrak c_F(H(t))\,dt\wedge d\bar t.
\tag{RC26}
\]

It is strictly positive away from \(t=0\) and has quadratic vanishing
there. Equivalently,

\[
\boxed{\left.\partial_t^2\partial_{\bar t}^2
\log\det H_F(t)\right|_{t=0}
=\frac{\mathfrak c_F(H(0))}{|u|^2}>0.}
\tag{RC27}
\]

The coefficient of \(t^2\bar t^2\) in \(\log\det H_F\) is therefore
\(\mathfrak c_F(H(0))/(4|u|^2)\); the factor four comes from \(2!2!\).
This mixed derivative is not the second derivative along the real axis.
For \(F=0\) or \(F=E\), the relevant curvature is identically zero;
(RC13) is essential for strict positivity.

\subsection{The first derivative still records the aggregate
defect}\label{the-first-derivative-still-records-the-aggregate-defect}

At \(t=0\) and real \(u>0\), (RC18) gives

\[
-u\left.\frac{d}{dt}
\log\frac{\det H_F(t)}{\det G_{N,F}}\right|_{t=0}-kp
=2\operatorname{Re}\operatorname{Tr}A_F-kp,
\qquad G_{N,F}=\iota^*G_N\iota.
\tag{RC28}
\]

This is the supplied constituent-minor observer, not a new estimate.
Taking \(F=E_>\), the complete positive-real-defect generalized
eigenspace, gives the existing aggregate spectral defect. Equations
(RC25)--(RC27) add the exact transverse curvature carried by that same
marked constituent. A derivative of the full determinant alone sees only
the balanced total trace on a reflection-stable packet.

\section{Transfer to the original theta
metric}\label{transfer-to-the-original-theta-metric}

\subsection{A source-space expression for the
coefficient}\label{a-source-space-expression-for-the-coefficient}

Using (RC2), the same coefficient for the actual arithmetic metric is

\[
\boxed{\mathfrak c_F(G_N)=
\left(\min_{f\in F}\|r_N(1_E-\iota f)\|^2\right)
\left(\sup_{0\ne f\in F}
\frac{|\ell_F(f)|^2}{\|r_N\iota f\|^2}\right).}
\tag{RC29}
\]

The supremum equals the finite dual norm in (RC22). Every source norm,
the arithmetic unit implicit in \(r_N\), and every allowed theta
correction remain in this expression.

The isometry from the original source metric into period coordinates is
not \(E\xrightarrow{\Pi}(\mathbb C^q,I_q)\). It is

\[
(E,G_N)\xrightarrow{\Pi(t)}
(\mathbb C^q,\widetilde G_N(t)),
\qquad \widetilde G_N(t)=\Pi(t)^{-*}G_N\Pi(t)^{-1}.
\tag{RC30}
\]

Consequently the corresponding quotient-unit/residue product is
invariant when the subspace, unit, functional, and metric are all
transported: \(F\mapsto\Pi F\), \(1_E\mapsto\Pi1_E\), and
\(\ell\mapsto\ell\Pi^{-1}\). In particular it equals
\(\mathfrak c_F(G_N)\), independently of \(t\). This does not identify
the Chern form for a varying target metric with the fixed-target-metric
form (RC26).

\subsection{A finite condition-number
bridge}\label{a-finite-condition-number-bridge}

Let \(a_N(t)>0\) and \(b_N(t)>0\) be the least and greatest eigenvalues
of

\[
B_N^{\rm per}(t)=G_N^{-1}H(t).
\tag{RC31}
\]

This operator is positive and self-adjoint in \(G_N\), so
\(a_NG_N\preceq H(t)\preceq b_NG_N\). Minimizing a quadratic form over a
fixed coset preserves this order, whereas inversion of the restricted
positive form reverses it. Therefore

\[
\frac{a_N}{b_N}\mathfrak c_F(G_N)
\leq\mathfrak c_F(H(t))
\leq\frac{b_N}{a_N}\mathfrak c_F(G_N).
\tag{RC32}
\]

Writing \(\mathcal K_N=b_N/a_N\), the exact curvature has the two-sided
source-sensitive enclosure

\[
\boxed{\frac{|t|^2}{|u|^2\mathcal K_N(t)}\mathfrak c_F(G_N)
\leq\partial_t\partial_{\bar t}\log\det H_F(t)
\leq\frac{|t|^2\mathcal K_N(t)}{|u|^2}\mathfrak c_F(G_N).}
\tag{RC33}
\]

The calculated determinant of \(B_N^{\rm per}\) determines its
eigenvalue product, not \(\mathcal K_N\). Equation (RC33) keeps the
precise additional metric quantity required to compare the two curvature
observations.

\subsection{The existing endpoint certificate controls variation of this
coefficient}\label{the-existing-endpoint-certificate-controls-variation-of-this-coefficient}

For the original nested source degrees \(i\leq j\), let
\(G_i\succeq G_j\succ0\), \(V_n=\det G_n\), and

\[
\theta=\lambda_{\min}(G_i^{-1}G_j)\in(0,1].
\tag{RC34}
\]

Then \(\theta G_i\preceq G_j\preceq G_i\). The quotient norms satisfy
\(\theta q_{G_i,F}\leq q_{G_j,F}\leq q_{G_i,F}\), and the squared
restricted dual norms satisfy
\(\|\ell_F\|_{G_i^{-1}}^2\leq\|\ell_F\|_{G_j^{-1}}^2
\leq\theta^{-1}\|\ell_F\|_{G_i^{-1}}^2\), with each inverse taken on the
restricted form. Multiplying gives

\[
\theta\mathfrak c_F(G_i)\leq\mathfrak c_F(G_j)
\leq\theta^{-1}\mathfrak c_F(G_i).
\tag{RC35}
\]

Every eigenvalue of \(G_i^{-1}G_j\) lies in \((0,1]\), and one equals
\(\theta\). Their product is at most \(\theta\). Hence

\[
\boxed{\left|\log\frac{\mathfrak c_F(G_j)}{\mathfrak c_F(G_i)}\right|
\leq-\log\theta\leq\log\frac{V_i}{V_j}.}
\tag{RC36}
\]

This is a new use of the existing endpoint quantity. Any certified upper
bound for the original \(\log(V_i/V_j)\) immediately bounds the change
of the specified residue/quotient coupling, without a separately assumed
uniform spectral gap. It does not reverse into an upper bound for
\(\log(V_i/V_j)\) from the coupling alone.

\section{The period realization of the arithmetic
Laplacian}\label{the-period-realization-of-the-arithmetic-laplacian}

The quadratic polynomial of the original arithmetic generator is

\[
\mathsf L_k=A^2-kA.
\tag{RC37}
\]

For \(k=1\) this is the polynomial corresponding to Connes--Consani's
\(H(1+H)\) under \(H=-D\). A nonconstant parameter connection introduces
an exact additional term at second order. From (RC5),

\[
\Pi''(t)=\frac1{u^2}\Pi(t)A(t)^2
-\frac1u\Pi(t)\mathcal R.
\tag{RC38}
\]

Thus, on every fibre,

\[
\boxed{u^2\Pi''(t)+ku\Pi'(t)+u\Pi(t)\mathcal R
=\Pi(t)(A(t)^2-kA(t)).}
\tag{RC39}
\]

At zero this is the exact period comparison for (RC37). Omitting
\(u\Pi\mathcal R\) would give a false equation. The correction is the
same residue map that supplies (RC15) and (RC25).

The original source lift of this correction is

\[
\mathcal R_N^{\rm src}=r_N\mathcal Rj_E:\mathscr X\to\mathscr X,
\qquad j_E\mathcal R_N^{\rm src}=\mathcal Rj_E,
\qquad \mathcal R_N^{\rm src}d=0.
\tag{RC40}
\]

It is defined on the specified source domain. For a fixed
\(x\in\mathscr X\) let \(Y_x(t)=\Pi(t)j_Ex\). Applying (RC39) at zero
and using (RC2) gives

\[
\begin{split}
\left.(u^2\partial_t^2+ku\partial_t)Y_x\right|_{0}
+u\Pi(0)j_E\mathcal R_N^{\rm src}x
=\Pi(0)j_E\bigl((D^{(k)})^2-kD^{(k)}\bigr)x.
\end{split}
\tag{RC41}
\]

The original primitive for \(D^{(k)}r_N-r_NA=dK_N\) is still present;
the perturbed source operator \(D^{(k)}+t\mathcal R_N^{\rm src}\) has
the same primitive, exactly as in DS69.

For the original Gram, retain

\[
W_N=A^*G_N+G_NA-kG_N.
\tag{RC42}
\]

Direct expansion proves

\[
\mathsf L_k^*G_N-G_N\mathsf L_k=A^*W_N-W_NA,
\qquad G_N\mathsf L_k=W_NA-A^*G_NA.
\tag{RC43}
\]

These are the preceding Laplacian-control identities, included here to
fix the map to the present period calculation. Under (RC30), put
\(\widetilde A=\Pi A\Pi^{-1}\),
\(\widetilde{\mathsf L}_k=\Pi\mathsf L_k\Pi^{-1}\), and
\(\widetilde W_N=\Pi^{-*}W_N\Pi^{-1}\). Substitution gives

\[
\widetilde A^*\widetilde G_N+
\widetilde G_N\widetilde A-k\widetilde G_N=\widetilde W_N,
\tag{RC44}
\]

and the two identities (RC43) in precisely the transported metric. Thus
every original numerical-range/control estimate is transported by an
actual isometry, not by exchanging \(G_N\) for \(H(t)\).

The full primary action is retained:

\[
\mathsf L_k|_{E_\rho}
=\rho(\rho-k)I+(2\rho-k)N_\rho+N_\rho^2.
\tag{RC45}
\]

The finite residue algebra and the adelic Hochschild source are not
being declared isomorphic. The latter reaches these finite coefficients
through the original theta complex, its quotient, and \(j_E\). In
particular a theta boundary is killed by \(j_E\); the differentiated
original relation producing the Jacobian trace is taken before that
quotient. The finite algebra's higher Hochschild modules from the
supplied periodic resolution, \(HH_{2a+1}=E/(\chi')\) and
\(HH_{2a+2}=\operatorname{ann}_E(\chi')\), remain separate from the
operator \(\mathcal R\) and from the adelic Hochschild groups.

\section{Translation, cohomological support, and exact
scope}\label{translation-cohomological-support-and-exact-scope}

\subsection{Translation preserves curvature but changes the marked
trace}\label{translation-preserves-curvature-but-changes-the-marked-trace}

The supplied DT1--13 comparison is

\[
\chi_a(S)=\chi(S-a),\quad
\Phi_a(S)=\Phi(S-a)-\Phi(-a),\quad
\Pi_a(t)=f_a(t)\Pi(t)C_a,
\quad f_a(t)=e^{(-\Phi(-a)-ta)/u},
\tag{RC46}
\]

where \(C_a\) is the Pascal substitution matrix, \(\det C_a=1\), and
\(C_aA_aC_a^{-1}=A+aI\). Transport the same constituent by
\(\iota_a=C_a^{-1}\iota\). Then

\[
P_{F,a}(t)=f_a(t)P_F(t),\qquad
\det H_{F,a}(t)=|f_a(t)|^{2p}\det H_F(t).
\tag{RC47}
\]

Since \(f_a\) is nowhere zero and holomorphic,

\[
\boxed{\partial_t\partial_{\bar t}\log\det H_{F,a}(t)
=\partial_t\partial_{\bar t}\log\det H_F(t).}
\tag{RC48}
\]

But, with the original \(k\) fixed,

\[
2\operatorname{Re}\operatorname{Tr}(A_F+aI)-kp
=(2\operatorname{Re}\operatorname{Tr}A_F-kp)+2p\operatorname{Re}a.
\tag{RC49}
\]

Thus curvature, unlike the first-derivative observer, does not retain
this scalar normalization by itself. The exact scalar \(f_a\) and the
marked arithmetic fibre \(a=0\) remain necessary parts of the
comparison. This is not an assertion that translated polynomials are
actual zeta packets. In the specified finite-field realization DT12--13
instead tensors with the constant Artin--Schreier character line, whose
geometric-Frobenius scalar is a root of unity. The finite-field weights
are preserved by that operation; it does not identify that Frobenius
with \(A\).

\subsection{Cohomological degrees and supported-zero
maps}\label{cohomological-degrees-and-supported-zero-maps}

For an invariant \(F\), the supplied logarithmic lattice
\(\mathcal M_F=F\mathcal O+E\mathcal O(-D)\) gives

\[
\mathbb H^0(\mathcal K_F)=F,\quad
\mathbb H^1(\mathcal K_F)=0,\quad
\mathbb H^2(\mathcal K_F)=(E/F)\otimes\mathbb T.
\tag{RC50}
\]

Its support-compatible dual is
\(\mathcal M_F^\vee(-D)=F^{\rm ann}\mathcal O+E^\vee\mathcal O(-D)\).
The evaluation target includes that divisor twist. The full coefficient
extension \(0\to F\to E\to E/F\to0\) is retained through the outer
comparison complexes; it is not reconstructed from the two middle
cohomology groups alone. These are the corrected distinctions in
DS25--34 and DS66--69.

The new transverse map (RC14) is therefore a specified map
\(\mathbb H^0(\mathcal K_F)\to E/F\), with the latter connected to
degree two by the retained line \(\mathbb T\). It is not silently a
degree-zero chain endomorphism of \(\mathcal K_F\). In particular,
\(\mathcal R\) need not preserve its lattice; (RC15) quantifies the
failure.

Lift every coefficient map by the original support reconstruction,
retaining the structural square

\[
\begin{array}{ccc}
G(\mathbb Z)&\longrightarrow&G(\mathbb C)\\
p_{\mathbb Z}\downarrow&&\downarrow p_{\mathbb C}\\
\mathbb Z&\longrightarrow&\mathbb C.
\end{array}
\tag{RC51}
\]

For each support label \(\lambda\), the new map is

\[
(\lambda,v)\longmapsto
(\lambda,(\pi1_E)\ell_F(v)),\qquad \tau\longmapsto\tau.
\tag{RC52}
\]

If \(v\in\ker\ell_F\), its receiving value is the supported zero at
\(\lambda\), not external absence. For the single active fibre,

\[
\mathsf S(f)^{-1}(\tau)=\{\tau\},\qquad
\mathsf S(f)^{-1}(e)=(\ker f)^\bullet,
\qquad p^{-1}(\ker f)=\{\tau\}\sqcup(\ker f)^\bullet.
\tag{RC53}
\]

This uses the corrected split-kernel formula in the supplied review. The
original infinite arithmetic quotient, support transitions, and all
source boundaries remain attached.

\section{Verification and provenance}\label{verification-and-provenance}

The included checker contains twelve independent finite test methods.
Its fixtures are monic polynomials with rational coefficients, including
repeated roots, and exact complex-rational matrices. It checks the
residue pairing and trace, the \(q^2\) endomorphism basis, the exact
transverse rank and kernel, the curvature factorization, a bivariate
fourth-order log-Gram jet with its factorial, the metric enclosures, the
forced Laplacian equation, and the constituent translation comparison.
The normal and optimized runs have their actual command and output
records in \texttt{checks/}. Three deliberately false formulas are
tested separately: zero transverse rank, omitted Laplacian forcing, and
omitted \(|t|^2\) in the curvature. These must fail in both modes.

The finite tests do not prove improper-integral convergence, the global
analytic period theorem, Deligne's theorem, or uniform growth on
arithmetic packets. Those parts have their written source proofs and
stated hypotheses. No predecessor checker or Lean build is claimed as a
new execution here. No remote repository or publication is modified by
this package.

The dated \emph{Deligne Primary Source Correction} supersedes the
historical unresolved-citation claim in \emph{Exponential Intake}: the
reference is Roman \textbf{(I.8.11)} to Weil I, printed page 306. The
strict weight inequality, the tensor square of the same eigenvalue, and
the lisse/mixed hypotheses are retained. This note relies on that
supplied correction for those printed-page identifications; it does not
claim a fresh full audit of Weil II.

The intake receipt records the hashes actually measured from the
uploaded files. All nine files listed in the manifest's \texttt{copies}
array match their output hashes. The currently supplied README is newer
than the README hash recorded in that manifest. The independent review
also pins translation TeX hash \texttt{39163624...}, whereas the
delivered translation source and the manifest agree on
\texttt{bd1831cd...}. The review's source-byte identity must therefore
be reconciled before its acceptance is attributed to those exact
delivered translation bytes; this discrepancy does not itself show that
a displayed mathematical formula is false.

\section{References and supplied-source
locators}\label{references-and-supplied-source-locators}

Bloch, S., \& Esnault, H. (1999). \emph{Gauß--Manin determinant
connections and periods for irregular connections} {[}Preprint{]}.
arXiv, math/9912095. The inspected §5 supplies the
polynomial-exponential period context, not the new assertions (RC11),
(RC15), (RC25), or (RC36).

Connes, A., \& Consani, C. (2023). Hochschild homology, trace map and
\(\zeta\)-cycles. In A. Connes, C. Consani, B. I. Dundas, M. Khalkhali,
\& H. Moscovici (Eds.), \emph{Cyclic cohomology at 40: Achievements and
future prospects} (Proceedings of Symposia in Pure Mathematics, Vol.
105, pp.~83--101). American Mathematical Society. DOI:
10.1090/pspum/105/01896.

Deligne, P. (1980). La conjecture de Weil: II. \emph{Publications
Mathématiques de l'IHÉS, 52}, 137--252. DOI: 10.1007/BF02684780. The
present note does not infer characteristic-zero spectral purity from
finite-field purity.

National Institute of Standards and Technology. (n.d.). §5.5: Functional
relations. In \emph{Digital Library of Mathematical Functions}. Equation
5.5.6, Gauss's gamma multiplication formula, gives the constant in
(RC6).

\emph{Deligne Split Cohomology}. (2026, September 13). Supplied research
sidebar, DS1--69. Source:
\texttt{deligne\_split\_sidebar\_20260913.tex}; relevant maps DS12--14,
DS25--34, DS53--69.

\emph{Exponential Period Determinant}. (2026, September 13). Supplied
research note, XD1--26. Source:
\texttt{deligne\_exponential\_determinant\_extension\_20260913.tex}.

\emph{Exponential Translation Comparison}. (2026, September 13).
Supplied research note, DT1--13. Source:
\texttt{Exponential\_Translation\_Comparison.tex}.

\emph{Pasted markdown (20260913-034115)}. (2026, September 13). Supplied
SGA trace/period continuation. Sections 2--3 give the residue
coevaluation and periodic resolution; §§4--5 evaluate and differentiate
the determinant; §§6--7 identify the constituent minor and transport the
original metric.

\emph{Independent Deligne Review}, \emph{Exponential Intake}, and
\emph{Deligne Primary Source Correction}. (2026, September 13). Supplied
review and correction records. Their historical execution and
source-reading claims are not relabeled as new executions in this note.

\end{document}
